\documentclass[11pt]{article}

\usepackage[a4paper,margin=1in]{geometry}
\usepackage{fontspec}
\setmainfont{Liberation Serif}
\setsansfont{Liberation Sans}
\setmonofont{DejaVu Sans Mono}[Scale=0.82]
\usepackage{amsmath,amssymb,amsthm,mathtools}
\usepackage{booktabs,array,longtable}
\usepackage{xcolor}
\usepackage{hyperref}
\usepackage{enumitem}
\usepackage{listings}
\usepackage{caption}

\hypersetup{colorlinks=true,linkcolor=blue,urlcolor=blue,citecolor=blue}

\definecolor{codebg}{RGB}{248,248,248}
\definecolor{codeframe}{RGB}{220,220,220}

\lstdefinestyle{leanstyle}{
  basicstyle=\ttfamily\scriptsize,
  breaklines=true,
  breakatwhitespace=false,
  columns=fullflexible,
  keepspaces=true,
  showstringspaces=false,
  frame=single,
  framerule=0.3pt,
  rulecolor=\color{codeframe},
  backgroundcolor=\color{codebg},
  xleftmargin=0.5em,
  xrightmargin=0.5em,
  aboveskip=0.8em,
  belowskip=0.8em
}
\lstset{style=leanstyle}

\setlist[itemize]{leftmargin=*}
\setlist[enumerate]{leftmargin=*}

\newcommand{\code}[1]{\texttt{#1}}
\newcommand{\status}[1]{\texttt{#1}}
\newcommand{\oblig}[1]{\textsf{#1}}

\title{TS20 -- Horizon Goldbach: A Certified Conditional Architecture\\
\large From Combinatorial Discharge to Analytic Infrastructure Ledger}
\author{Serge Durand\\Horizon / Goldbach Lean Programme}
\date{Baseline: Lean 4.15.0 / Mathlib v4.15.0}

\begin{document}
\maketitle

\begin{abstract}
We present the final synthesis of the Horizon Goldbach formal specification
programme (TS14--TS19). The programme decomposes a conditional route to the
binary Goldbach conjecture into a modular Lean 4 architecture. Four layers are
established: (1) a conditional gluing theorem (TS14); (2) a localisation of the
analytic residue into the interface \code{ShortIntervalPrimeSecondMoment}
(TS15); (3) a purely combinatorial discharge (TS16) that eliminates the finite
counting obligation; and (4) three relative discharge theorems (TS17, TS18,
TS19) that reduce the remaining analytic obstacles to named local
infrastructures. In the audited TS15--TS19 specification layer, no \code{sorry}
and no global \code{axiom} remain. The final ledger records exactly what
analytic work is needed to turn the conditional architecture into an
unconditional proof. This document serves as a roadmap for future formalisation
efforts and as an invitation to the analytic number theory community.
\end{abstract}

\section{Introduction}

The Horizon programme develops a Lean 4 formal specification of a conditional
proof architecture for the binary Goldbach conjecture, following the general
shape of the Hardy--Littlewood circle method. The central TS14 layer is a
conditional collage:
\[
  h_R \wedge h_G \wedge h_S \wedge h_C
  \Longrightarrow
  \forall N\ge 4,\; 2\mid N,\; \mathrm{GoldbachAt}(N),
\]
where \(h_R\) is the Riemann Hypothesis, \(h_G\) is a strong large sieve input,
\(h_S\) is a spectral renormalisation bound, and \(h_C\) is the computational
verification up to \(4\cdot10^{18}\) in the style of Oliveira e Silva et al.

The purpose of TS14--TS19 is not to prove these analytic and computational
inputs. It is to expose their exact role, remove avoidable finite or
combinatorial debt, and isolate the remaining analytic content as local,
typed, auditable infrastructure.

\section{Summary of the Sprints}

\subsection{TS14 -- Conditional collage}

TS14 states the conditional theorem and records the macroscopic hypotheses. The
proof splits the even integers into a finite interval, covered by computation,
and an asymptotic interval, covered by analytic estimates. This is the global
conditional layer.

\subsection{TS15 -- Localisation of the analytic residue}

TS15 introduces the interface \code{ShortIntervalPrimeSecondMoment}, encoding a
short-interval second-moment estimate for primes. It shows that the downstream
problem \code{Problem_E1} follows from this interface together with a
finite-counting comparison between prime-pair counts and short-interval
energy.

\subsection{TS16 -- Combinatorial discharge}

TS16 proves the finite-counting comparison in full generality. The Lean theorem
\code{pair_count_le_energy} injects close pairs into energetic triples and
uses only finite set/cardinality reasoning. This part is
\status{repo_committed}: no analytic assumption is involved.

\subsection{TS17 -- Mellin--Jackson discharge}

TS17 replaces the abstract Mellin--Jackson projection interface by a relative
theorem. It proves that the TS15 projection bound follows from two local
infrastructures:
\begin{itemize}
  \item \oblig{MellinFourierNormBridge}, for the logarithmic change of
    variables and the identification of \(\Theta\) with \(d/du\);
  \item \oblig{FourierTailInfrastructure}, for the Plancherel-based
    high-frequency tail estimate.
\end{itemize}
The result is \status{repo_committed_relative}.

\subsection{TS18 -- Short-interval second moment discharge}

TS18 reduces \code{ShortIntervalPrimeSecondMoment} to two local structures:
\begin{itemize}
  \item \oblig{DirichletCharacterBridge}, for character orthogonality,
    discrete Parseval, and bridge errors;
  \item \oblig{LargeSieveInfrastructure}, for a local large-sieve estimate
    producing a constant \(C\le 1\).
\end{itemize}
The relative instance \code{secondMomentInstance} and the downstream theorem
\code{Problem_E1_from_TS18} are proved mechanically. The result is
\status{repo_committed_relative}.

\subsection{TS19 -- OTSA residual discharge}

TS19 decomposes the OTSA residual estimate. It introduces:
\begin{itemize}
  \item \oblig{KernelSpectralControl}, with constant \(C_k\);
  \item \oblig{TraceContributionControl}, with constant \(C_t\);
  \item \oblig{MellinTailDecay}, with constant \(C_m\);
  \item \oblig{OTSACouplingHypothesis}, the local coupling estimate.
\end{itemize}
The residual estimate is reduced to
\[
  R(x,B) \le (C_k C_t + C_m) S(x,B),
\]
and Lean proves that if \(C_k C_t + C_m \le 26\), then
\[
  R(x,B) \le 26\, S(x,B).
\]
The constant \(26\) is therefore an explicit threshold condition.

\section{Final Analytic Infrastructure Ledger}

The audited specification layer leaves the following named obligations. They
can be counted either as eight individual structures or as six analytic
families, grouping the three OTSA constants and their coupling as one OTSA
block.

\begin{center}
\small
\begin{longtable}{p{0.30\linewidth}p{0.55\linewidth}}
\toprule
Obligation & Description \\
\midrule
\oblig{MellinFourierNormBridge} &
Links the Mellin tail norm to the Fourier tail norm via \(u=\log x\) and
identifies \(\Theta\) with \(d/du\). \\
\oblig{FourierTailInfrastructure} &
Provides the Plancherel-based tail estimate
\(\|(I-P_T)B\|_2 \le T^{-k}\|B^{(k)}\|_2\). \\
\oblig{DirichletCharacterBridge} &
Relates the short-interval prime energy to a character-side second moment,
plus controlled bridge error. \\
\oblig{LargeSieveInfrastructure} &
Supplies a local large-sieve bound with constant \(C\le 1\) for the relevant
moduli and normalisation. \\
\oblig{KernelSpectralControl} &
Gives the OTSA spectral-kernel constant \(C_k\). \\
\oblig{TraceContributionControl} &
Gives the OTSA trace/pole constant \(C_t\). \\
\oblig{MellinTailDecay} &
Gives the OTSA Mellin-tail constant \(C_m\). \\
\oblig{OTSACouplingHypothesis} &
States the local inequality \(R(x,B)\le (C_kC_t+C_m)S(x,B)\). \\
\bottomrule
\end{longtable}
\end{center}

Once these obligations are instantiated by genuine analytic proofs in Lean, the
relative layers TS17--TS19 become unconditional. Together with the TS16
combinatorial discharge and the TS14 conditional gluing layer, this supplies a
fully formal route to the Goldbach conclusion under the corresponding TS14
global hypotheses.

\section{Audit of the Specification Layers}

The audited scope is:
\[
  \code{TS/Goldbach/Strong/TS15}
  \quad\text{through}\quad
  \code{TS/Goldbach/Strong/TS19}.
\]
In this scope, the Lean code contains no \code{sorry} and no global
\code{axiom}. The analytic assumptions are passed as explicit structures and
parameters.

A global audit can be run as:

\begin{lstlisting}[language={}]
rg -n "s[o]rry" TS/Goldbach/Strong/TS15 TS/Goldbach/Strong/TS16 \
                 TS/Goldbach/Strong/TS17 TS/Goldbach/Strong/TS18 \
                 TS/Goldbach/Strong/TS19
rg -n "a[x]iom" TS/Goldbach/Strong/TS15 TS/Goldbach/Strong/TS16 \
                 TS/Goldbach/Strong/TS17 TS/Goldbach/Strong/TS18 \
                 TS/Goldbach/Strong/TS19
\end{lstlisting}

Expected output: no matches.

\section{Build Commands}

The core Lean targets can be built with:

\begin{lstlisting}[language={}]
lake build TS.Goldbach.Strong.TS16.CombinatorialDischarge \
          TS.Goldbach.Strong.TS17.MellinJacksonDischarge \
          TS.Goldbach.Strong.TS18.SecondMomentDischarge \
          TS.Goldbach.Strong.TS19.OTSAResidualDischarge
\end{lstlisting}

The local baseline used for the audited files is Lean 4.15.0.

\section{Conclusion}

The Horizon programme has constructed a machine-auditable reduction of the
Goldbach proof architecture to a compact set of named analytic infrastructure
obligations. The finite combinatorial core is proven unconditionally. The
remaining analytic estimates are isolated, exposed, and can be worked on
independently.

The architecture is modular, auditable, and ready to receive future
contributions, whether they are formalisations of classical theorems
(Mellin--Fourier bridge, large sieve) or targeted estimates for the OTSA
residual block. The present document, together with the Lean 4 source files of
TS15--TS19, specifies exactly what remains to be done.

\appendix

\section{Lean Source Index}

The associated source tree contains the following audited sprint directories:

\begin{itemize}
  \item \code{TS/Goldbach/Strong/TS15}
  \item \code{TS/Goldbach/Strong/TS16}
  \item \code{TS/Goldbach/Strong/TS17}
  \item \code{TS/Goldbach/Strong/TS18}
  \item \code{TS/Goldbach/Strong/TS19}
\end{itemize}

TS20 is a synthesis document and does not add new Lean theorems.

\end{document}
